\chapter{Joint Densities, Transformations and Order Statistics}
\label{ch:jointcont}

We now treat several continuous random variables together. The key ideas are the
joint density, the change-of-variables formula in several dimensions (the
Jacobian), and the resulting formulas for sums, ratios and order statistics. We
finish with the multivariate Gaussian distribution, the most important joint law
in the subject.

\section{Joint densities}

\begin{definition}[Joint density]
Random variables $X,Y$ have \emph{joint density} $f_{X,Y}:\R^2\to[0,\infty)$ if
for every (reasonable) set $A\subseteq\R^2$,
\[
\Prob\bigl((X,Y)\in A\bigr)=\iint_Af_{X,Y}(x,y)\,dx\,dy.
\]
\end{definition}

The marginal densities are obtained by integrating out the other variable:
\[
f_X(x)=\int_{-\infty}^\infty f_{X,Y}(x,y)\,dy,\qquad
f_Y(y)=\int_{-\infty}^\infty f_{X,Y}(x,y)\,dx .
\]
(Reason: $\Prob(X\le x)=\int_{-\infty}^x\int_{-\infty}^\infty f_{X,Y}(s,y)\,dy\,ds$,
which exhibits the inner integral as a density for $X$.) The LOTUS formula
$\Exp[g(X,Y)]=\iint g\,f_{X,Y}\,dx\,dy$ is proved as in
Theorem~\ref{thm:lotuscont}.

\begin{definition}[Independence]
Continuous random variables $X,Y$ are \emph{independent} if
$\Prob(X\in A,Y\in B)=\Prob(X\in A)\Prob(Y\in B)$ for all intervals $A,B$.
\end{definition}

\begin{theorem}[Independence and factorisation]\label{thm:indfactor}
If $(X,Y)$ has a joint density, then $X,Y$ are independent iff
$f_{X,Y}(x,y)=f_X(x)f_Y(y)$ for (almost) all $(x,y)$. More generally, if the joint
density can be written $f_{X,Y}(x,y)=g(x)h(y)$ for some functions $g,h\ge0$, then
$X,Y$ are independent and $f_X\propto g$, $f_Y\propto h$.
\end{theorem}
\begin{proof}
If $f_{X,Y}=f_Xf_Y$ then $\Prob(X\in A,Y\in B)=\int_A\int_Bf_X(x)f_Y(y)\,dy\,dx=
\Prob(X\in A)\Prob(Y\in B)$. For the converse, independence gives
$\iint_{A\times B}f_{X,Y}=\iint_{A\times B}f_Xf_Y$ for all rectangles, and two
non-negative functions with equal integrals over all rectangles agree almost
everywhere (justified in Chapter~\ref{ch:measure}). For the last claim, the
marginals are $f_X=g\cdot\int h$ and $f_Y=h\cdot\int g$, and
$1=\iint gh=\int g\int h$, so $f_Xf_Y=gh\,(\int g\int h)=gh$.
\end{proof}

\begin{nuance}{the support must be a product set}
The factorisation $f(x,y)=g(x)h(y)$ is only possible if the support is a
rectangle (or a product of sets). If $(X,Y)$ is uniform on the triangle
$\{0<y<x<1\}$ with density $2$ there, then $X,Y$ are not independent: knowing
$Y=0.9$ forces $X>0.9$. The indicator $\indic_{\{y<x\}}$ cannot be factored.
\end{nuance}

\begin{example}[Uniform on a disc]
Let $(X,Y)$ be uniform on the unit disc $D$: $f=1/\pi$ on $D$. The marginal is
$f_X(x)=\int_{-\sqrt{1-x^2}}^{\sqrt{1-x^2}}\frac1\pi\,dy=\frac2\pi\sqrt{1-x^2}$ for
$\abs x<1$, a semicircle law, \emph{not} uniform. $X$ and $Y$ are uncorrelated (by
symmetry $\Exp[XY]=0$) but dependent, since the support is not a rectangle.
\end{example}

\begin{example}[Meeting problem]
Two people arrive independently and uniformly in a one-hour window. They meet if
they arrive within $\delta$ hours of each other ($0<\delta<1$). With
$(X,Y)$ uniform on $[0,1]^2$, the event $\{\abs{X-Y}<\delta\}$ has complement two
corner triangles of legs $1-\delta$, so
$\Prob(\text{meet})=1-(1-\delta)^2$. For $\delta=1/4$ this is $7/16$.
\end{example}

\section{Sums and convolution}

\begin{theorem}[Convolution formula]\label{thm:convcont}
If $X,Y$ are independent with densities $f_X,f_Y$, then $X+Y$ has density
\[
f_{X+Y}(z)=\int_{-\infty}^\infty f_X(x)\,f_Y(z-x)\,dx .
\]
\end{theorem}
\begin{proof}
$\Prob(X+Y\le z)=\iint_{x+y\le z}f_X(x)f_Y(y)\,dy\,dx=\int f_X(x)F_Y(z-x)\,dx$.
Differentiating under the integral in $z$ gives the formula.
\end{proof}

\begin{theorem}[Sum of independent gammas]\label{thm:gammasum}
If $X\sim\mathrm{Gamma}(\alpha,\lambda)$ and $Y\sim\mathrm{Gamma}(\beta,\lambda)$ are
independent then $X+Y\sim\mathrm{Gamma}(\alpha+\beta,\lambda)$. In particular, a sum of
$n$ independent $\mathrm{Exp}(\lambda)$ variables is $\mathrm{Gamma}(n,\lambda)$.
\end{theorem}
\begin{proof}
For $z>0$,
\[
f_{X+Y}(z)=\frac{\lambda^{\alpha+\beta}}{\Gamma(\alpha)\Gamma(\beta)}e^{-\lambda z}
\int_0^zx^{\alpha-1}(z-x)^{\beta-1}\,dx
=\frac{\lambda^{\alpha+\beta}}{\Gamma(\alpha)\Gamma(\beta)}z^{\alpha+\beta-1}e^{-\lambda z}
\,B(\alpha,\beta),
\]
by substituting $x=zt$. This is proportional to the Gamma$(\alpha+\beta,\lambda)$
density; since both integrate to $1$, the constants agree (so we also learn
$B(\alpha,\beta)=\Gamma(\alpha)\Gamma(\beta)/\Gamma(\alpha+\beta)$, derived again
below).
\end{proof}

\begin{theorem}[Sum of independent normals]\label{thm:normalsum}
If $X\sim N(\mu_1,\sigma_1^2)$ and $Y\sim N(\mu_2,\sigma_2^2)$ are independent then
$X+Y\sim N(\mu_1+\mu_2,\sigma_1^2+\sigma_2^2)$.
\end{theorem}
\begin{proof}
It suffices to consider means zero (shift) and use the convolution formula. The
integrand has exponent $-\tfrac12\bigl[x^2/\sigma_1^2+(z-x)^2/\sigma_2^2\bigr]$. Complete
the square in $x$: with $s^2=\sigma_1^2\sigma_2^2/(\sigma_1^2+\sigma_2^2)$ and
$m=z\sigma_1^2/(\sigma_1^2+\sigma_2^2)$,
\[
\frac{x^2}{\sigma_1^2}+\frac{(z-x)^2}{\sigma_2^2}
=\frac{(x-m)^2}{s^2}+\frac{z^2}{\sigma_1^2+\sigma_2^2}.
\]
(Verify by expanding both sides.) The $x$-integral of $\exp(-(x-m)^2/(2s^2))$ is
$s\sqrt{2\pi}$, independent of $z$, so
$f_{X+Y}(z)\propto\exp\bigl(-z^2/(2(\sigma_1^2+\sigma_2^2))\bigr)$, the
$N(0,\sigma_1^2+\sigma_2^2)$ density.
\end{proof}

\section{Change of variables in several dimensions}

\begin{theorem}[Multivariate change of variables]\label{thm:jacobian}
Let $(X,Y)$ have joint density $f_{X,Y}$ concentrated on an open set $D$, and let
$T:D\to D'$, $(x,y)\mapsto(u,v)=(g_1(x,y),g_2(x,y))$ be a continuously
differentiable bijection with continuously differentiable inverse
$(u,v)\mapsto(x,y)=(h_1(u,v),h_2(u,v))$. Then $(U,V)=T(X,Y)$ has joint density
\[
f_{U,V}(u,v)=f_{X,Y}\bigl(h_1(u,v),h_2(u,v)\bigr)\,\abs{J(u,v)},\qquad
J=\det\begin{pmatrix}\partial x/\partial u&\partial x/\partial v\\
\partial y/\partial u&\partial y/\partial v\end{pmatrix}.
\]
\end{theorem}
\begin{proof}
For any set $B\subseteq D'$, $\Prob((U,V)\in B)=\Prob((X,Y)\in T^{-1}(B))=
\iint_{T^{-1}(B)}f_{X,Y}(x,y)\,dx\,dy$. The change of variables theorem for
multiple integrals (from calculus) with $(x,y)=(h_1,h_2)(u,v)$ converts this to
$\iint_Bf_{X,Y}(h_1,h_2)\abs J\,du\,dv$. Since this holds for all $B$, the
integrand is the density of $(U,V)$.
\end{proof}

\begin{example}[Gamma--Beta relation]\label{ex:gammabeta}
Let $X\sim\mathrm{Gamma}(a,1)$ and $Y\sim\mathrm{Gamma}(b,1)$ be independent, and put
$U=X+Y$, $V=X/(X+Y)$. The inverse map is $x=uv$, $y=u(1-v)$, with Jacobian
\[
J=\det\begin{pmatrix}v&u\\1-v&-u\end{pmatrix}=-uv-u(1-v)=-u,\qquad\abs J=u.
\]
So for $u>0$, $0<v<1$,
\begin{align*}
f_{U,V}(u,v)&=\frac{(uv)^{a-1}e^{-uv}}{\Gamma(a)}\cdot\frac{(u(1-v))^{b-1}e^{-u(1-v)}}{\Gamma(b)}\cdot u\\
&=\underbrace{\frac{u^{a+b-1}e^{-u}}{\Gamma(a)\Gamma(b)}}_{\text{function of }u}\;
\underbrace{v^{a-1}(1-v)^{b-1}}_{\text{function of }v}.
\end{align*}
By Theorem~\ref{thm:indfactor}, $U$ and $V$ are \emph{independent}. Integrating out
$u$ gives the marginal of $V$: $\frac{\Gamma(a+b)}{\Gamma(a)\Gamma(b)}v^{a-1}(1-v)^{b-1}$, with
$U\sim\mathrm{Gamma}(a+b,1)$. Since this must integrate to $1$,
\[
B(a,b)=\frac{\Gamma(a)\Gamma(b)}{\Gamma(a+b)},\qquad V\sim\mathrm{Beta}(a,b).
\]
Consequently $\Exp[V]=B(a+1,b)/B(a,b)=\frac{\Gamma(a+1)\Gamma(a+b)}{\Gamma(a)\Gamma(a+b+1)}
=\frac a{a+b}$, as quoted in Chapter~\ref{ch:cont}.
\end{example}

\begin{example}[Box--Muller transformation]\label{ex:boxmuller}
Let $X,Y$ be independent $N(0,1)$ and let $(R,\Theta)$ be polar coordinates:
$x=r\cos\theta$, $y=r\sin\theta$, $\abs J=r$. Then
\[
f_{R,\Theta}(r,\theta)=\frac1{2\pi}e^{-r^2/2}\cdot r,\qquad r>0,\ 0<\theta<2\pi .
\]
This factorises: $\Theta\sim U(0,2\pi)$ and $R$ has density $re^{-r^2/2}$,
independent of $\Theta$. Further, $R^2$ has CDF
$\Prob(R^2\le s)=\int_0^{\sqrt s}re^{-r^2/2}dr=1-e^{-s/2}$, so $R^2\sim\mathrm{Exp}(1/2)$.
Reading this backwards gives a way to generate normals: if $U_1,U_2\sim U(0,1)$
are independent, then $R=\sqrt{-2\ln U_1}$ and $\Theta=2\pi U_2$ give
$X=R\cos\Theta$, $Y=R\sin\Theta$ independent standard normals.
\end{example}

\begin{example}[Ratio of normals is Cauchy]
Let $X,Y$ be independent $N(0,1)$ and $W=Y/X$. In polar coordinates
$W=\tan\Theta$, and $\Theta$ is uniform on $(0,2\pi)$ by the previous example, so
$\Prob(W\le w)=\Prob(\tan\Theta\le w)=\frac12+\frac1\pi\arctan w$, which is the
Cauchy CDF.
\end{example}

\section{Order statistics}

Let $X_1,\dots,X_n$ be independent with common density $f$ and CDF $F$. Sorting
them gives the \emph{order statistics}
$X_{(1)}\le X_{(2)}\le\dots\le X_{(n)}$ (ties have probability zero). In
particular $X_{(1)}=\min$ and $X_{(n)}=\max$.

\begin{theorem}[Distribution of the $k$th order statistic]\label{thm:orderstat}
\[
\Prob(X_{(k)}\le x)=\sum_{j=k}^n\binom nj F(x)^j\bigl(1-F(x)\bigr)^{n-j},
\qquad
f_{X_{(k)}}(x)=\frac{n!}{(k-1)!\,(n-k)!}F(x)^{k-1}\bigl(1-F(x)\bigr)^{n-k}f(x).
\]
\end{theorem}
\begin{proof}
$X_{(k)}\le x$ iff at least $k$ of the $X_i$ are $\le x$. The number of $X_i\le x$ is
Bin$(n,F(x))$, which gives the CDF. For the density, differentiate term by term. Put
$a_j=\binom{n-1}jF^j(1-F)^{n-1-j}$ (with $a_n=0$). Using
$j\binom nj=n\binom{n-1}{j-1}$ and $(n-j)\binom nj=n\binom{n-1}{j}$,
\[
\frac d{dx}\Bigl[\binom njF^j(1-F)^{n-j}\Bigr]
=nf\,\bigl(a_{j-1}-a_j\bigr).
\]
Summing over $j=k,\dots,n$ telescopes to $nf\,a_{k-1}$, and
$n\binom{n-1}{k-1}=n!/((k-1)!(n-k)!)$.
\end{proof}

\begin{example}[Uniform order statistics]\label{ex:unifos}
For $U_i\sim U(0,1)$, $F(x)=x$, so
$f_{U_{(k)}}(x)=\frac{n!}{(k-1)!(n-k)!}x^{k-1}(1-x)^{n-k}$, that is,
$U_{(k)}\sim\mathrm{Beta}(k,n-k+1)$. Therefore
\[
\Exp[U_{(k)}]=\frac k{n+1},\qquad\Var(U_{(k)})=\frac{k(n-k+1)}{(n+1)^2(n+2)} .
\]
In particular $\Exp[\max]=n/(n+1)$ and $\Exp[\min]=1/(n+1)$. The $n$ points divide
$[0,1]$ into $n+1$ gaps, and by symmetry each gap has mean length $1/(n+1)$.
\end{example}

\begin{theorem}[Minimum of independent exponentials]\label{thm:minexp}
If $X_i\sim\mathrm{Exp}(\lambda_i)$ are independent, then
$\min_iX_i\sim\mathrm{Exp}(\sum_i\lambda_i)$, and the index of the minimum is $i$
with probability $\lambda_i/\sum_j\lambda_j$.
\end{theorem}
\begin{proof}
$\Prob(\min>t)=\prod_i\Prob(X_i>t)=e^{-(\sum\lambda_i)t}$. For the index, by
conditioning on $X_i=t$ and using independence (LOTUS):
$\Prob(X_i\text{ is min})=\int_0^\infty\lambda_ie^{-\lambda_it}\prod_{j\ne i}e^{-\lambda_jt}dt=
\frac{\lambda_i}{\sum_j\lambda_j}$.
\end{proof}

\begin{interview}{expected max/min}
For $n$ i.i.d. $U(0,1)$: $\Exp[\max]=\frac n{n+1}$ and $\Exp[\min]=\frac1{n+1}$. For
$n$ i.i.d. $\mathrm{Exp}(1)$: $\Exp[\min]=\frac1n$ and
$\Exp[\max]=H_n=\sum_{k=1}^n\frac1k$ (the time until the last of $n$ lightbulbs
fails; by memorylessness the successive gaps are independent exponentials with
rates $n,n-1,\dots,1$). These results recur constantly.
\end{interview}

\section{The multivariate normal distribution}

\begin{definition}[Standard and general multivariate normal]
A random vector $Z=(Z_1,\dots,Z_d)$ is \emph{standard multivariate normal} if its
components are independent $N(0,1)$ variables. A vector $X$ is \emph{multivariate
normal} (jointly Gaussian) if $X=\mu+AZ$ for some $d\times m$ matrix $A$, vector
$\mu\in\R^d$, and standard normal $Z\in\R^m$. Its covariance matrix is
$\Sigma=AA^{\mathsf T}$.
\end{definition}

Every linear combination $\sum c_iX_i=c^{\mathsf T}\mu+(A^{\mathsf T}c)^{\mathsf T}Z$
is a sum of independent normals, hence univariate normal (Theorem~\ref{thm:normalsum}).
This is the characterising property: $X$ is jointly Gaussian iff every linear
combination of its components is normal.

\begin{theorem}[Density]\label{thm:mvndensity}
If $\Sigma$ is invertible (equivalently $A$ is square and invertible), $X=\mu+AZ$ has
density
\[
f_X(x)=\frac1{(2\pi)^{d/2}\sqrt{\det\Sigma}}
\exp\Bigl(-\tfrac12(x-\mu)^{\mathsf T}\Sigma^{-1}(x-\mu)\Bigr).
\]
\end{theorem}
\begin{proof}
$Z$ has density $(2\pi)^{-d/2}e^{-\abs z^2/2}$. The map $x=\mu+Az$ has inverse
$z=A^{-1}(x-\mu)$ with Jacobian determinant $\det A^{-1}$, and
$\abs{\det A}=\sqrt{\det(AA^{\mathsf T})}=\sqrt{\det\Sigma}$. Also
$\abs z^2=(x-\mu)^{\mathsf T}(A^{-1})^{\mathsf T}A^{-1}(x-\mu)=(x-\mu)^{\mathsf T}\Sigma^{-1}(x-\mu)$.
Apply Theorem~\ref{thm:jacobian} (its $d$-dimensional version).
\end{proof}

\subsection*{The bivariate case}

For $d=2$ take $X=Z_1$, $Y=\rho Z_1+\sqrt{1-\rho^2}\,Z_2$ with $\abs\rho<1$. Then
$X,Y\sim N(0,1)$, $\Cov(X,Y)=\rho$, and the density is
\[
f(x,y)=\frac1{2\pi\sqrt{1-\rho^2}}\exp\Bigl(-\frac{x^2-2\rho xy+y^2}{2(1-\rho^2)}\Bigr).
\]
\emph{Derivation.} The inverse map is $z_1=x$, $z_2=(y-\rho x)/\sqrt{1-\rho^2}$, with
Jacobian $1/\sqrt{1-\rho^2}$, and
$z_1^2+z_2^2=x^2+(y-\rho x)^2/(1-\rho^2)=(x^2-2\rho xy+y^2)/(1-\rho^2)$.

\begin{theorem}[Properties of the bivariate normal]\label{thm:bvn}
\begin{enumerate}[label=(\roman*)]
\item If $(X,Y)$ is jointly Gaussian, then $X,Y$ are independent iff they are
uncorrelated.
\item The conditional law of $Y$ given $X=x$ (the density $f(x,y)/f_X(x)$ as a
function of $y$) is $N(\rho x,\,1-\rho^2)$.
\end{enumerate}
\end{theorem}
\begin{proof}
(i) Independent implies uncorrelated always. Conversely if $\rho=0$ the density above
factorises into $\varphi(x)\varphi(y)$. (In general, uncorrelated jointly Gaussian
variables have a block-diagonal $\Sigma$, and the density factorises.)

(ii) Dividing, and using $f_X=\varphi$:
\[
\frac{f(x,y)}{\varphi(x)}=\frac1{\sqrt{2\pi(1-\rho^2)}}
\exp\Bigl(-\frac{x^2-2\rho xy+y^2}{2(1-\rho^2)}+\frac{x^2}2\Bigr)
=\frac1{\sqrt{2\pi(1-\rho^2)}}\exp\Bigl(-\frac{(y-\rho x)^2}{2(1-\rho^2)}\Bigr),
\]
since $-(x^2-2\rho xy+y^2)+(1-\rho^2)x^2=-(y-\rho x)^2$.
\end{proof}

\begin{pitfall}{normal marginals do not imply joint normality}
Let $X\sim N(0,1)$, and let $S=\pm1$ with probability $\tfrac12$ each, independent
of $X$. Put $Y=SX$. Then $Y\sim N(0,1)$ (by symmetry), and $\Cov(X,Y)=\Exp[S]\Exp[X^2]=0$, so
$X,Y$ are uncorrelated and each is normal. But they are \emph{not} independent
(since $\abs Y=\abs X$), hence not jointly Gaussian. The statement
``uncorrelated implies independent'' needs \emph{joint} normality.
\end{pitfall}

\begin{nuance}{conditioning is linear for Gaussians}
Part (ii) shows the best predictor of $Y$ from $X$ is linear in $X$:
$\Exp[Y\given X]=\rho X$, and the conditional variance $1-\rho^2$ does not depend on
$x$. This is special to Gaussians, and it is the reason linear regression and the
Kalman filter work so cleanly.
\end{nuance}

\section*{Exercises}
\addcontentsline{toc}{section}{Exercises}

\begin{enumerate}
\item $(X,Y)$ has joint density $f(x,y)=x+y$ on $[0,1]^2$. Find the marginals,
$\Cov(X,Y)$, and decide whether $X,Y$ are independent.
\item Let $X,Y$ be independent $U(0,1)$. Find the density of $X+Y$ (triangular) and
$\Prob(X+Y>1.5)$.
\item Let $X\sim\mathrm{Exp}(\lambda)$, $Y\sim\mathrm{Exp}(\mu)$ independent. Find
$\Prob(X<Y)$ and the density of $X-Y$.
\item Show that if $X,Y$ are independent $\mathrm{Exp}(1)$ then $X/(X+Y)\sim U(0,1)$.
\item Show that the Cauchy distribution is stable: if $X_1,X_2$ are independent
standard Cauchy then $(X_1+X_2)/2$ is standard Cauchy. (Use the convolution
formula; this is laborious, and a much shorter proof appears in
Chapter~\ref{ch:transforms}.)
\item Find $\Exp[\max(X_1,X_2,X_3)]$ for i.i.d.\ $N(0,1)$ variables. (Use
symmetry: $\Exp[\max]=-\Exp[\min]$, and the sum of the order statistics is
known.)
\item $n$ points are placed independently and uniformly on a circle of
circumference $1$. Find the expected length of the arc containing a given point.
\item Let $(X,Y)$ be standard bivariate normal with correlation $\rho$. Compute
$\Prob(X>0,Y>0)$. (Hint: $\Theta$-symmetry of the rotation-invariant case, then
the general case via $\frac14+\frac{\arcsin\rho}{2\pi}$.)
\item (Challenge) Derive the joint density of $(X_{(1)},X_{(n)})$ for i.i.d.\ continuous
random variables, and find $\Exp[X_{(n)}-X_{(1)}]$ (the expected range) for
uniform samples.
\end{enumerate}
